\documentclass[10pt,a4paper]{amsart}
\usepackage{amsmath,amssymb,amsthm,graphicx}
% Retained preamble for the Stacks Project chapter snapshot in this directory.
% The chapter's own preamble.tex is a copy of the shared Stacks Project preamble; the block
% below is sliced MECHANICALLY, line for line, from it. Nothing was added, removed or
% redefined; the wrapper class, the xr-hyper/hyperref wiring and the externaldocument list
% are omitted because the assembler supplies the document class.
%
% lines taken from preamble.tex: 3, 5, 9, 11, 12, 156, 157, 158, ... (51 lines total)
%
% package coverage restored from the chapter's own preamble (2026-10-13): xy, xy-2cell

\usepackage[all]{xy}\usepackage{verbatim}
\xyoption{2cell}
\UseAllTwocells

\newenvironment{reference}{\comment}{\endcomment}
\newenvironment{slogan}{\comment}{\endcomment}
\newenvironment{history}{\comment}{\endcomment}
\theoremstyle{plain}
\newtheorem{theorem}[subsection]{Theorem}
\newtheorem{proposition}[subsection]{Proposition}
\newtheorem{lemma}[subsection]{Lemma}
\theoremstyle{definition}
\newtheorem{definition}[subsection]{Definition}
\newtheorem{example}[subsection]{Example}
\newtheorem{exercise}[subsection]{Exercise}
\newtheorem{situation}[subsection]{Situation}
\theoremstyle{remark}
\newtheorem{remark}[subsection]{Remark}
\newtheorem{remarks}[subsection]{Remarks}
\numberwithin{equation}{subsection}
\def\lim{\mathop{\mathrm{lim}}\nolimits}
\def\colim{\mathop{\mathrm{colim}}\nolimits}
\def\Spec{\mathop{\mathrm{Spec}}}
\def\Hom{\mathop{\mathrm{Hom}}\nolimits}
\def\Ext{\mathop{\mathrm{Ext}}\nolimits}
\def\SheafHom{\mathop{\mathcal{H}\!\mathit{om}}\nolimits}
\def\SheafExt{\mathop{\mathcal{E}\!\mathit{xt}}\nolimits}
\def\Sch{\mathit{Sch}}
\def\Mor{\mathop{\mathrm{Mor}}\nolimits}
\def\Ob{\mathop{\mathrm{Ob}}\nolimits}
\def\Sh{\mathop{\mathit{Sh}}\nolimits}
\def\NL{\mathop{N\!L}\nolimits}
\def\CH{\mathop{\mathrm{CH}}\nolimits}
\def\proetale{{pro\text{-}\acute{e}tale}}
\def\etale{{\acute{e}tale}}
\def\QCoh{\mathit{QCoh}}
\def\Ker{\mathop{\mathrm{Ker}}}
\def\Im{\mathop{\mathrm{Im}}}
\def\Coker{\mathop{\mathrm{Coker}}}
\def\Coim{\mathop{\mathrm{Coim}}}
\DeclareMathSymbol{\boxtimes}{\mathbin}{AMSa}{"02}
\def\QCohstack{\mathcal{QC}\!\mathit{oh}}
\def\Cohstack{\mathcal{C}\!\mathit{oh}}
\def\Spacesstack{\mathcal{S}\!\mathit{paces}}
\def\Quotfunctor{\mathrm{Quot}}
\def\Hilbfunctor{\mathrm{Hilb}}
\def\Curvesstack{\mathcal{C}\!\mathit{urves}}
\def\Polarizedstack{\mathcal{P}\!\mathit{olarized}}
\def\Complexesstack{\mathcal{C}\!\mathit{omplexes}}
\def\Pic{\mathop{\mathrm{Pic}}\nolimits}
\def\Picardstack{\mathcal{P}\!\mathit{ic}}
\def\Picardfunctor{\mathrm{Pic}}
\def\Deformationcategory{\mathcal{D}\!\mathit{ef}}
\usepackage{cleveref}
\crefname{originalproposition}{original proposition}{original propositions}
\Crefname{originalproposition}{Original proposition}{Original propositions}
\makeatletter
\newlabel{definition-elementary-distinguished-square}{{08GM}{source}{}{sourcepointer-1}{}}
\newlabel{definition-elementary-distinguished-square@cref}{{[originalproposition][0][]08GM}{[1][0][]source}{}{}{}}
\makeatother
\crefname{originalproposition}{original proposition}{original propositions}
\Crefname{originalproposition}{Original proposition}{Original propositions}
\makeatletter
\newlabel{derived-remark-check-on-generator}{{0ATH}{source}{}{sourcepointer-2}{}}
\newlabel{derived-remark-check-on-generator@cref}{{[originalproposition][0][]0ATH}{[1][0][]source}{}{}{}}
\makeatother
\crefname{originalproposition}{original proposition}{original propositions}
\Crefname{originalproposition}{Original proposition}{Original propositions}
\makeatletter
\newlabel{lemma-bounded-truncation}{{0GFF}{source}{}{sourcepointer-3}{}}
\newlabel{lemma-bounded-truncation@cref}{{[originalproposition][0][]0GFF}{[1][0][]source}{}{}{}}
\makeatother
\crefname{originalproposition}{original proposition}{original propositions}
\Crefname{originalproposition}{Original proposition}{Original propositions}
\makeatletter
\newlabel{lemma-direct-summand-of-a-restriction}{{09IX}{source}{}{sourcepointer-4}{}}
\newlabel{lemma-direct-summand-of-a-restriction@cref}{{[originalproposition][0][]09IX}{[1][0][]source}{}{}{}}
\makeatother
\crefname{originalproposition}{original proposition}{original propositions}
\Crefname{originalproposition}{Original proposition}{Original propositions}
\makeatletter
\newlabel{lemma-ext-from-perfect-into-bounded-QCoh}{{09MB}{source}{}{sourcepointer-5}{}}
\newlabel{lemma-ext-from-perfect-into-bounded-QCoh@cref}{{[originalproposition][0][]09MB}{[1][0][]source}{}{}{}}
\makeatother
\crefname{originalproposition}{original proposition}{original propositions}
\Crefname{originalproposition}{Original proposition}{Original propositions}
\makeatletter
\newlabel{lemma-induction-principle}{{08GP}{source}{}{sourcepointer-6}{}}
\newlabel{lemma-induction-principle@cref}{{[originalproposition][0][]08GP}{[1][0][]source}{}{}{}}
\makeatother
\crefname{originalproposition}{original proposition}{original propositions}
\Crefname{originalproposition}{Original proposition}{Original propositions}
\makeatletter
\newlabel{lemma-pushforward-perfect}{{08HK}{source}{}{sourcepointer-7}{}}
\newlabel{lemma-pushforward-perfect@cref}{{[originalproposition][0][]08HK}{[1][0][]source}{}{}{}}
\makeatother
\crefname{originalproposition}{original proposition}{original propositions}
\Crefname{originalproposition}{Original proposition}{Original propositions}
\makeatletter
\newlabel{lemma-pushforward-with-support-in-open}{{08GG}{source}{}{sourcepointer-8}{}}
\newlabel{lemma-pushforward-with-support-in-open@cref}{{[originalproposition][0][]08GG}{[1][0][]source}{}{}{}}
\makeatother
\crefname{originalproposition}{original proposition}{original propositions}
\Crefname{originalproposition}{Original proposition}{Original propositions}
\makeatletter
\newlabel{lemma-quasi-coherence-direct-image}{{08FA}{source}{}{sourcepointer-9}{}}
\newlabel{lemma-quasi-coherence-direct-image@cref}{{[originalproposition][0][]08FA}{[1][0][]source}{}{}{}}
\makeatother
\crefname{originalproposition}{original proposition}{original propositions}
\Crefname{originalproposition}{Original proposition}{Original propositions}
\makeatletter
\newlabel{perfect-lemma-orthogonal-koszul-first-variant}{{0GEJ}{source}{}{sourcepointer-10}{}}
\newlabel{perfect-lemma-orthogonal-koszul-first-variant@cref}{{[originalproposition][0][]0GEJ}{[1][0][]source}{}{}{}}
\makeatother
\crefname{originalproposition}{original proposition}{original propositions}
\Crefname{originalproposition}{Original proposition}{Original propositions}
\makeatletter
\newlabel{perfect-proposition-detecting-bounded-above}{{0GEN}{source}{}{sourcepointer-11}{}}
\newlabel{perfect-proposition-detecting-bounded-above@cref}{{[originalproposition][0][]0GEN}{[1][0][]source}{}{}{}}
\makeatother
\crefname{originalproposition}{original proposition}{original propositions}
\Crefname{originalproposition}{Original proposition}{Original propositions}
\makeatletter
\newlabel{remark-classical-generator}{{0GFC}{source}{}{sourcepointer-12}{}}
\newlabel{remark-classical-generator@cref}{{[originalproposition][0][]0GFC}{[1][0][]source}{}{}{}}
\makeatother
\crefname{originalproposition}{original proposition}{original propositions}
\Crefname{originalproposition}{Original proposition}{Original propositions}
\makeatletter
\newlabel{sites-cohomology-lemma-dual-perfect-complex}{{08JJ}{source}{}{sourcepointer-13}{}}
\newlabel{sites-cohomology-lemma-dual-perfect-complex@cref}{{[originalproposition][0][]08JJ}{[1][0][]source}{}{}{}}
\makeatother
\crefname{originalproposition}{original proposition}{original propositions}
\Crefname{originalproposition}{Original proposition}{Original propositions}
\makeatletter
\newlabel{sites-cohomology-section-global-RHom}{{0B6E}{source}{}{sourcepointer-14}{}}
\newlabel{sites-cohomology-section-global-RHom@cref}{{[originalproposition][0][]0B6E}{[1][0][]source}{}{}{}}
\makeatother
\crefname{originalproposition}{original proposition}{original propositions}
\Crefname{originalproposition}{Original proposition}{Original propositions}
\makeatletter
\newlabel{theorem-bondal-van-den-Bergh}{{09IY}{source}{}{sourcepointer-15}{}}
\newlabel{theorem-bondal-van-den-Bergh@cref}{{[originalproposition][0][]09IY}{[1][0][]source}{}{}{}}
\makeatother
\title{Stacks, Derived Categories of Spaces: detecting-bounded-above}
\author{The Stacks Project Authors}
\date{stacks-project:spaces-perfect.tex:local-snapshot-20261012; GNU FDL 1.2 or later}
\begin{document}
\maketitle
\noindent\textit{Editorial note.} Statement span is lines 7298-7322 of sources/stacks-spaces-perfect/pinned-original-source.tex; proof span is lines 7324-7419, the complete author proof (94 body lines).\par\medskip
\noindent\textit{Editorial note.} Navigation stubs: 15 cited-but-not-included labels are supplied as non-load-bearing pointers with their pinned tags where the snapshot tag index has them.\par\medskip
\noindent\textit{Editorial note.} extra\_packages is empty on purpose: the retained preamble already loads xy with options and a second declaration clashes in the assembler wrapper.\par\medskip
\noindent\textit{Editorial note.} Text is retained verbatim, including the author's own typos and placeholders.\par\medskip
\noindent\textit{Editorial note.} Queue index 53 of queue\_spaceperfect.json (pb=4363, stmt 7298-7322, proof 7324-7419).\par\medskip
\noindent\textit{Editorial note.} stacks-spaces-perfect wave sp01 (2026-10-14): fresh adjudication accept per lane-s3/sp01\_judgments; pool per Chief 12:4x unfreeze (reserve-pool workflow, six-check clean); spaces-perfect.tex 126 with-proof rows, 70 b107-occupied labels avoided, 56 free rows judged 32 accept / 24 reject.\par\medskip
\noindent\textit{Editorial note.} Cross-chapter same-label warning (NOT an occupancy conflict, Chief ruling 2026-10-13 evening): label 'proposition-detecting-bounded-above' also exists in stacks-perfect — a different proposition in a different chapter file.\par\medskip
\noindent\textit{Allowed standard prerequisites (editorial).}\begin{itemize}
\item Explicit prerequisite boundary. The cited inputs are prerequisites; the argument in the retained span is not replaced by any cited result.
\end{itemize}\medskip
\ifdefined\hypertarget\hypertarget{sourcepointer-1}{}\fi
\noindent\textit{Source reference.} Original proposition 08GM: sources/stacks-spaces-perfect/pinned-original-source.tex, source lines 1513--1528. Cited by the retained statement or proof but not part of the collected spans; supplied only as a navigation pointer. The argument does not depend on its content.\par\medskip
\ifdefined\hypertarget\hypertarget{sourcepointer-2}{}\fi
\noindent\textit{Source reference.} Original proposition 0ATH: sources/stacks-spaces-perfect/pinned-original-source.tex, source lines 7419--7419. Cited by the retained statement or proof but not part of the collected spans; supplied only as a navigation pointer. The argument does not depend on its content.\par\medskip
\ifdefined\hypertarget\hypertarget{sourcepointer-3}{}\fi
\noindent\textit{Source reference.} Original proposition 0GFF: sources/stacks-spaces-perfect/pinned-original-source.tex, source lines 7253--7261. Cited by the retained statement or proof but not part of the collected spans; supplied only as a navigation pointer. The argument does not depend on its content.\par\medskip
\ifdefined\hypertarget\hypertarget{sourcepointer-4}{}\fi
\noindent\textit{Source reference.} Original proposition 09IX: sources/stacks-spaces-perfect/pinned-original-source.tex, source lines 3877--3884. Cited by the retained statement or proof but not part of the collected spans; supplied only as a navigation pointer. The argument does not depend on its content.\par\medskip
\ifdefined\hypertarget\hypertarget{sourcepointer-5}{}\fi
\noindent\textit{Source reference.} Original proposition 09MB: sources/stacks-spaces-perfect/pinned-original-source.tex, source lines 4328--4338. Cited by the retained statement or proof but not part of the collected spans; supplied only as a navigation pointer. The argument does not depend on its content.\par\medskip
\ifdefined\hypertarget\hypertarget{sourcepointer-6}{}\fi
\noindent\textit{Source reference.} Original proposition 08GP: sources/stacks-spaces-perfect/pinned-original-source.tex, source lines 1563--1577. Cited by the retained statement or proof but not part of the collected spans; supplied only as a navigation pointer. The argument does not depend on its content.\par\medskip
\ifdefined\hypertarget\hypertarget{sourcepointer-7}{}\fi
\noindent\textit{Source reference.} Original proposition 08HK: sources/stacks-spaces-perfect/pinned-original-source.tex, source lines 3610--3616. Cited by the retained statement or proof but not part of the collected spans; supplied only as a navigation pointer. The argument does not depend on its content.\par\medskip
\ifdefined\hypertarget\hypertarget{sourcepointer-8}{}\fi
\noindent\textit{Source reference.} Original proposition 08GG: sources/stacks-spaces-perfect/pinned-original-source.tex, source lines 2135--2145. Cited by the retained statement or proof but not part of the collected spans; supplied only as a navigation pointer. The argument does not depend on its content.\par\medskip
\ifdefined\hypertarget\hypertarget{sourcepointer-9}{}\fi
\noindent\textit{Source reference.} Original proposition 08FA: sources/stacks-spaces-perfect/pinned-original-source.tex, source lines 790--804. Cited by the retained statement or proof but not part of the collected spans; supplied only as a navigation pointer. The argument does not depend on its content.\par\medskip
\ifdefined\hypertarget\hypertarget{sourcepointer-10}{}\fi
\noindent\textit{Source reference.} Original proposition 0GEJ: sources/stacks-spaces-perfect/pinned-original-source.tex, source lines 7419--7419. Cited by the retained statement or proof but not part of the collected spans; supplied only as a navigation pointer. The argument does not depend on its content.\par\medskip
\ifdefined\hypertarget\hypertarget{sourcepointer-11}{}\fi
\noindent\textit{Source reference.} Original proposition 0GEN: sources/stacks-spaces-perfect/pinned-original-source.tex, source lines 7419--7419. Cited by the retained statement or proof but not part of the collected spans; supplied only as a navigation pointer. The argument does not depend on its content.\par\medskip
\ifdefined\hypertarget\hypertarget{sourcepointer-12}{}\fi
\noindent\textit{Source reference.} Original proposition 0GFC: sources/stacks-spaces-perfect/pinned-original-source.tex, source lines 4141--4153. Cited by the retained statement or proof but not part of the collected spans; supplied only as a navigation pointer. The argument does not depend on its content.\par\medskip
\ifdefined\hypertarget\hypertarget{sourcepointer-13}{}\fi
\noindent\textit{Source reference.} Original proposition 08JJ: sources/stacks-spaces-perfect/pinned-original-source.tex, source lines 7419--7419. Cited by the retained statement or proof but not part of the collected spans; supplied only as a navigation pointer. The argument does not depend on its content.\par\medskip
\ifdefined\hypertarget\hypertarget{sourcepointer-14}{}\fi
\noindent\textit{Source reference.} Original proposition 0B6E: sources/stacks-spaces-perfect/pinned-original-source.tex, source lines 7419--7419. Cited by the retained statement or proof but not part of the collected spans; supplied only as a navigation pointer. The argument does not depend on its content.\par\medskip
\ifdefined\hypertarget\hypertarget{sourcepointer-15}{}\fi
\noindent\textit{Source reference.} Original proposition 09IY: sources/stacks-spaces-perfect/pinned-original-source.tex, source lines 3891--3901. Cited by the retained statement or proof but not part of the collected spans; supplied only as a navigation pointer. The argument does not depend on its content.\par\medskip
\begin{proposition}
\label{proposition-detecting-bounded-above}
Let $S$ be a scheme.
Let $X$ be a quasi-compact and quasi-separated algebraic space over $S$.
Let $G \in D_{perf}(\mathcal{O}_X)$ be a perfect complex which generates
$D_\QCoh (\mathcal{O}_X)$. Let $E \in D_\QCoh (\mathcal{O}_X)$.
The following are equivalent
\begin{enumerate}
\item $E \in D^-_\QCoh (\mathcal{O}_X)$,
\item $\Hom_{D(\mathcal{O}_X)}(G[-i], E) = 0$ for $i \gg 0$,
\item $\Ext^i_X(G, E) = 0$ for $i \gg 0$,
\item $R\Hom_X(G, E)$ is in $D^-(\mathbf{Z})$,
\item $H^i(X, G^\vee \otimes_{\mathcal{O}_X}^\mathbf{L} E) = 0$
for $i \gg 0$,
\item $R\Gamma(X, G^\vee \otimes_{\mathcal{O}_X}^\mathbf{L} E)$
is in $D^-(\mathbf{Z})$,
\item for every perfect object $P$ of $D(\mathcal{O}_X)$
\begin{enumerate}
\item the assertions (2), (3), (4) hold with $G$ replaced by $P$, and
\item $H^i(X, P \otimes_{\mathcal{O}_X}^\mathbf{L} E) = 0$ for $i \gg 0$,
\item $R\Gamma(X, P \otimes_{\mathcal{O}_X}^\mathbf{L} E)$
is in $D^-(\mathbf{Z})$.
\end{enumerate}
\end{enumerate}
\end{proposition}

\begin{proof}
Assume (1). Since
$\Hom_{D(\mathcal{O}_X)}(G[-i], E) = \Hom_{D(\mathcal{O}_X)}(G, E[i])$
we see that this is zero for $i \gg 0$ by
Lemma \ref{lemma-ext-from-perfect-into-bounded-QCoh}. This proves
that (1) implies (2).

\medskip\noindent
Parts (2), (3), (4) are equivalent by the discussion in
Cohomology on Sites, Section \ref{sites-cohomology-section-global-RHom}.
Part (5) and (6) are equivalent as $H^i(X, -) = H^i(R\Gamma(X, -))$
by definition. The equivalent conditions (2), (3), (4) are
equivalent to the equivalent conditions (5), (6) by
Cohomology on Sites, Lemma \ref{sites-cohomology-lemma-dual-perfect-complex}
and the fact that $(G[-i])^\vee = G^\vee[i]$.

\medskip\noindent
It is clear that (7) implies (2). Conversely,
let us prove that the equivalent conditions (2) -- (6) imply (7).
Recall that $G$ is a  classical generator for $D_{perf}(\mathcal{O}_X)$ by
Remark \ref{remark-classical-generator}.
For $P \in D_{perf}(\mathcal{O}_X)$ let $T(P)$ be the assertion that
$R\Hom_X(P, E)$ is in $D^-(\mathbf{Z})$.
Clearly, $T$ is inherited by direct sums,
satisfies the 2-out-of-three property for distinguished
triangles, is inherited by direct summands, and is preserved by shifts.
Hence by Derived Categories, Remark \ref{derived-remark-check-on-generator}
we see that (4) implies $T$ holds on all of $D_{perf}(\mathcal{O}_X)$.
The same argument works for all other properties, except that for property
(7)(b) and (7)(c) we also use that $P \mapsto P^\vee$ is a self
equivalence of $D_{perf}(\mathcal{O}_X)$. Small detail omitted.

\medskip\noindent
We will prove the equivalent conditions (2) -- (7) imply (1)
using the induction principle of
Lemma \ref{lemma-induction-principle}.

\medskip\noindent
First, we prove (2) -- (7) $\Rightarrow$ (1) if $X$ is affine.
This follows from the case of schemes, see
Derived Categories of Schemes, Proposition
\ref{perfect-proposition-detecting-bounded-above}.

\medskip\noindent
Now assume $(U \subset X, j : V \to X)$ is an elementary distinguished
square of quasi-compact and quasi-separated algebraic spaces over $S$
and assume the implication (2) -- (7) $\Rightarrow$ (1)
is known for $U$, $V$, and $U \times_X V$. To finish the proof
we have to show the implication (2) -- (7) $\Rightarrow$ (1) for $X$.
Suppose $E \in D_\QCoh(\mathcal{O}_X)$ satisfies (2) -- (7).
By Lemma \ref{lemma-direct-summand-of-a-restriction} and
Theorem \ref{theorem-bondal-van-den-Bergh} there exists a perfect
complex $Q$ on $X$ such that $Q|_U$ generates $D_\QCoh (\mathcal{O}_U)$.

\medskip\noindent
Say $V = \Spec(A)$. Let $Z \subset V$ be the reduced closed subscheme
which is the inverse image of $X \setminus U$ and maps isomorphically to it
(see Definition \ref{definition-elementary-distinguished-square}).
This is a retrocompact closed subset of $V$.
Choose $f_1, \ldots, f_r \in A$ such that
$Z = V(f_1, \ldots, f_r)$. Let $K \in D(\mathcal{O}_V)$ be the perfect
object corresponding to the Koszul complex on $f_1, \ldots, f_r$ over $A$.
Note that since $K$ is supported on $Z$, the pushforward
$K' = Rj_*K$ is a perfect object of $D(\mathcal{O}_X)$ whose
restriction to $V$ is $K$ (see Lemmas \ref{lemma-pushforward-perfect}
and \ref{lemma-pushforward-with-support-in-open}).
By assumption, we know $R\Hom_{\mathcal{O}_X}(Q, E)$ and
$R\Hom_{\mathcal{O}_X}(K', E)$ are bounded above.

\medskip\noindent
By Lemma \ref{lemma-pushforward-with-support-in-open}
we have $K' =  j_!K$ and hence
$$
\Hom_{D(\mathcal{O}_X)}(K'[-i], E) = \Hom_{D(\mathcal{O}_V)}(K[-i], E|_V) = 0
$$
for $i \gg 0$. Therefore, we may apply
Derived Categories of Schemes,
Lemma \ref{perfect-lemma-orthogonal-koszul-first-variant} to $E|_V$ to
obtain an integer $a$ such that
$\tau_{\geq a}(E|_V) =
\tau_{\geq a} R (U \times_X V \to V)_* (E|_{U \times_X V})$.
Then $\tau_{\geq a} E = \tau_{\geq a} R (U \to X)_* (E |_U)$
(check that the canonical map is an isomorphism after restricting to
$U$ and to $V$). Hence using Lemma \ref{lemma-bounded-truncation}
twice we see that
$$
\Hom_{D(\mathcal{O}_U)}(Q|_U [-i], E|_U) =
\Hom_{D(\mathcal{O}_X)}(Q[-i], R (U \to X)_* (E|_U)) = 0
$$
for $i \gg 0$. Since the Proposition holds for $U$ and the generator
$Q|_U$, we have $E|_U \in D^-_\QCoh(\mathcal{O}_U)$. But then since
the functor $R (U \to X)_*$ preserves $D^-_\QCoh$
(by Lemma \ref{lemma-quasi-coherence-direct-image}), we get
$\tau_{\geq a}E \in D^-_\QCoh(\mathcal{O}_X)$. Thus
$E \in D^-_\QCoh (\mathcal{O}_X)$.
\end{proof}

\end{document}
